\documentclass[12pt]{article}
\usepackage[utf8]{inputenc}
\usepackage[T1]{fontenc}
\usepackage[margin=1in]{geometry}
\usepackage{setspace}
\usepackage{xcolor}
\usepackage[colorlinks=true, urlcolor=blue!60!black, linkcolor=blue!60!black]{hyperref}
\setlength{\parindent}{2em}
\setlength{\parskip}{1.25em}
\renewcommand{\baselinestretch}{1.0}
\begin{document}
\begin{center}
{\LARGE \textit{Craig v. Missouri}, 29 U.S. (4 Pet.) 410 (1830)\par}
\vspace{0.8em}
{\large Case Brief --- Four-Page Edition\par}
\end{center}
\vspace{1em}

\section{Memory Jogger}

For a 4--3 Court, Chief Justice Marshall held that Missouri's interest-bearing loan-office certificates, issued on the faith of the state and designed to circulate as money, were unconstitutional ``bills of credit'' under Article I, Section 10, even though they were not legal tender, so a note given for them was void (\href{https://supreme.justia.com/cases/federal/us/29/410/}{29 U.S. 410}).

\section{Detailed Case Facts}

After the Panic of 1819 and amid a shortage of specie on the frontier, Missouri enacted a loan-office act in 1821 authorizing up to \$200,000 in paper certificates in denominations from fifty cents to ten dollars, bearing 2\% interest. The certificates were receivable for taxes and debts due the state, payable for officers' salaries and fees, and receivable for salt from the state's salt springs. Commissioners lent them to citizens on mortgages or personal security. The act pledged a redemption fund and declared that ``the faith of the state is hereby also pledged for the same purpose,'' with one-tenth withdrawn each year. The certificates were not legal tender for private debts but were designed to circulate as a local medium of exchange.

On August 1, 1822, Hiram Craig, John Moore, and Ephraim Moore gave the State a promissory note for \$199.99 plus interest at the Chariton loan office in exchange for certificates.

\section{Procedural History}

When they defaulted, Missouri sued in assumpsit. The defendants argued that the certificates were unconstitutional bills of credit, so the note lacked lawful consideration. The circuit court ruled for Missouri, and the Supreme Court of Missouri affirmed. On writ of error, with Senator Thomas Hart Benton arguing for Missouri, the U.S. Supreme Court reversed on March 12, 1830.

\section{Judicial Votes}

Reversed and remanded, \textbf{4--3}.
\begin{itemize}
\item \textbf{Majority:} Chief Justice Marshall, joined by Justices Duvall, Story, and Baldwin.
\item \textbf{Concurring:} None.
\item \textbf{Dissenting:} Justices Johnson, Thompson, and McLean, each separately.
\end{itemize}

\section{Holding}

(1) Section 25 jurisdiction existed because the record showed that the statute's constitutionality was necessarily decided below. (2) The certificates were bills of credit: paper issued by the state, on the faith of the state, intended to circulate as money, whatever their label and despite lacking legal-tender status. (3) The note given for them rested on illegal consideration and was void.

\section{Analysis of Opinions}

\textbf{Majority (Marshall).} \emph{Jurisdiction:} a state court cannot shield a federal question from review by omitting a recital; under \textit{Martin v.\ Hunter's Lessee}, the record showed that the only defense was the act's constitutionality. \emph{Definition:} read in light of colonial and Revolutionary usage, the clause does not bar ordinary state borrowing but forbids state paper meant to circulate as money. Missouri's certificates had every hallmark: issued by state officers in small denominations, lent into circulation, receivable for taxes, salaries, and salt, and backed by the state's faith. \emph{Substance over form:} the Constitution prohibits things, not labels. \emph{Independence from tender:} Article I, Section 10 separately forbids emitting bills of credit and making anything but gold and silver a tender; reading tender into bills of credit would erase a clause, and colonial non-tender notes, such as Virginia's, were bills of credit too. \emph{Remedy:} a note founded on an unconstitutional act cannot be enforced.

\textbf{Baldwin's vote.} In \textit{Briscoe} (1837), Baldwin revealed that he supplied the fourth vote on narrower grounds, believing a bill of credit required both a pledge of state faith and tender status for some debts.

\textbf{Johnson (dissenting).} Interest made the certificates' value change daily, unfit for currency, and they were promises to receive, not to pay.

\textbf{Thompson (dissenting).} A bill of credit rests merely on the drawer's credit; Missouri pledged specific funds and compelled no one to accept the certificates.

\textbf{McLean (dissenting).} A bill of credit requires compelled circulation without a redemption fund, and federal and state powers intermix ``like the colours of the rainbow,'' counseling ``mutual forbearance.''

\textbf{Later developments.} In \textit{Briscoe v.\ Bank of Commonwealth of Kentucky}, McLean wrote for a 6--1 Court upholding notes of a state-owned but separately incorporated bank, over Story's dissent. \textit{Poindexter v.\ Greenhow} and \textit{Houston \& Texas Central R.R.} narrowed the clause further.

\section{Examples: Future Applications}

\begin{enumerate}
\item \textbf{Same-Side (Invalid): Fiscal Liquidity Certificates.} A State prints small bearer certificates receivable for taxes and backed by its full faith and credit. Despite lacking tender status, they are bills of credit.
\item \textbf{Same-Side (Invalid): Interest-Bearing Scrip.} Scrip paid to employees, backed by lottery revenue, raises the very arguments the Johnson and Thompson dissents lost.
\item \textbf{Opposite-Side (Valid): Investment Bonds.} Large general-obligation bonds sold to investors are ordinary borrowing, not a circulating medium.
\item \textbf{Opposite-Side (Valid under \textit{Briscoe}): A State-Owned Bank.} Notes issued by a suable, separately incorporated state bank on its own capital are not bills of credit.
\item \textbf{Fence-Sitter: A State Digital Token.} A state corporation issues fully reserved, redeemable digital dollars disclaiming state faith. \textit{Briscoe} suggests validity, but \textit{Craig}'s things-not-names principle points the other way.
\end{enumerate}

\section{Critique}

\textbf{Formalism and functionalism.} Read with \textit{Briscoe}, \textit{Craig} exposes a classic tension: Marshall insisted that the Constitution bars things, not names, yet \textit{Briscoe} let states issue nearly identical currency through a state-owned corporation, raising the question whether Article I, Section 10's limits survive corporate form.

\textbf{Monetary history and windfall.} Economic historians and the dissenters note that Missouri's capped, interest-bearing, non-tender certificates were a modest response to specie scarcity, not the forced-tender paper money of 1786 Rhode Island; Marshall's rule favored national monetary discipline over state experimentation. And because the remedy voided the note, Craig and the Moores kept the value of the certificates and escaped repayment, a recurring dilemma when structural limits are enforced defensively in private debt litigation.

\section{Key Quotations}

\begin{quote}
``To emit bills of credit conveys to the mind the idea of issuing paper, intended to circulate through the community for its ordinary purposes as money, which paper is redeemable at a future day; this is the sense in which the terms have been always understood.''
\end{quote}
\noindent --- Chief Justice Marshall, Opinion of the Court, 29 U.S. at 432.

\begin{quote}
``Is the proposition to be maintained that the Constitution meant to prohibit names, and not things? That a very important act, big with great and ruinous mischief, which is expressly forbidden by words most appropriate for its description may be performed by the substitution of a name? That the Constitution, in one of its most important provisions, may be openly evaded by giving a new name to an old thing? We cannot think so.''
\end{quote}
\noindent --- Chief Justice Marshall, Opinion of the Court, 29 U.S. at 433.

\begin{quote}
``The Constitution therefore considers the emission of bills of credit and the enactment of tender laws as distinct operations, independent of each other, which may be separately performed; both are forbidden. To sustain the one because it is not also the other---to say, that bills of credit may be emitted if they be not made a tender in payment of debts---is in effect to expunge that distinct, independent prohibition and to read the clause as if it had been entirely omitted. We are not at liberty to do this.''
\end{quote}
\noindent --- Chief Justice Marshall, Opinion of the Court, 29 U.S. at 434.

\end{document}
